% =============================================================================
% VARIÁVEIS INSTRUMENTAIS (16:25 a 16:37)   [Script 06: exemplo completo]
% =============================================================================
\section[IV]{Variáveis instrumentais}

\begin{frame}{A ideia: um empurrão que não passa pela motivação}
  \centering
  \begin{tikzpicture}[scale=0.82,every node/.style={transform shape}]
    \node[cxi,minimum width=2.4cm] (z) at (0,0) {\textbf{Instrumento $Z$}\\sorteio da vaga};
    \node[cx,minimum width=2.4cm] (d) at (4,0) {\textbf{Participa $D$}\\fez o curso};
    \node[cx,minimum width=2.4cm] (y) at (8,0) {\textbf{Resultado $Y$}\\renda};
    \node[cxl,minimum width=2.8cm] (u) at (6,2.0) {\textbf{Motivação $U$}\\não observada};
    \draw[seta] (z) -- node[above,font=\scriptsize]{1) empurra} (d);
    \draw[seta] (d) -- node[above,font=\scriptsize]{efeito?} (y);
    \draw[setal] (u) -- (d); \draw[setal] (u) -- (y);
    \draw[seta,dashed,red] (z) to[bend right=35] node[below,font=\scriptsize]{2) sem caminho direto} (y);
  \end{tikzpicture}
  \vspace{0.2cm}
  \begin{itemize}
    \item Participar é escolha, e quem escolhe é mais motivado (\alert{endogeneidade});
    \item O \textbf{sorteio de vagas} muda quem participa sem depender da motivação. Efeito $=$ efeito da oferta $\div$ aumento da participação.
  \end{itemize}
\end{frame}

\begin{frame}{Variáveis instrumentais: características e exemplos}
  \carac{variação na participação \textbf{induzida pelo instrumento}}
        {\textbf{relevância} (move a participação); \textbf{exogeneidade} (como aleatório); \textbf{exclusão} (só afeta $Y$ via $D$)}
        {efeito local para os \textbf{\emph{compliers}} (LATE)}
        {um instrumento válido e forte (primeiro estágio com $F>10$; Staiger e Stock, 1997)}
        {adesão imperfeita; participação endógena com sorteio, regra ou distância}
  \vspace{0.05cm}
  \textbf{\footnotesize\color{azul}Exemplos em políticas públicas}\par
    \scriptsize
    \begin{itemize}
      \item \textbf{EUA:} sorteio para se inscrever no Medicaid em Oregon, para medir o efeito do seguro-saúde (Finkelstein et al., 2012).
      \item \textbf{Colômbia:} sorteio de vouchers escolares, para o efeito de usar o voucher (Angrist et al., 2002).
      \item \textbf{EUA:} loteria do alistamento militar e renda de veteranos (Angrist, 1990).
      \item \textbf{Hipotético:} sorteio de vagas do Qualifica (script 06, com o passo a passo completo).
    \end{itemize}
\end{frame}

\begin{frame}{IV em gráficos (script 06)}
  \begin{columns}[T]
    \begin{column}{0.42\textwidth}
      \centering
      \includegraphics[width=\textwidth]{06_iv_primeiro_estagio.png}\par
      \tiny \textbf{Primeiro estágio:} a oferta eleva a participação de 38\% para 75\% (+36 p.p.; $F=767$).
    \end{column}
    \begin{column}{0.58\textwidth}
      \centering
      \includegraphics[width=\textwidth]{06_iv_resumo.png}\par
      \tiny Oferta (ITT) $=107$; IV $=107/0{,}36\approx297$ (verdade: 300). Ingênua: 525. Se a oferta afetasse a renda por outro caminho, o IV iria a 574.
    \end{column}
  \end{columns}
\end{frame}
